\chapter{Gauge fields, curvature, and the meaning of a complex}
\label{chap:gauge}

\section{The precise connection to the preceding algebra}

The theory of orders concerns modules carrying a multiplication action.  The
theory of complexes concerns maps whose consecutive compositions vanish.  Gauge
theory adds a third ingredient: the maps are differential operators depending on
a connection, and the obstruction to their composing to zero is curvature.
This observation supplies a concrete bridge between the subjects, but it also
identifies the hypothesis that any proposed bridge must preserve.  An arbitrary
covariant exterior derivative is not a differential in a cochain complex.
Its square is a curvature action.  A flat connection removes that obstruction;
an anti-self-dual connection removes the appropriate projected obstruction.
Neither operation turns a classical deformation problem into a construction of
the quantum vacuum.

The following derivations fix conventions sufficiently explicitly to permit
comparison with both algebraic complexes and formal proofs.  The foundational
gauge theory dates to Yang and Mills; the geometric analysis of the
anti-self-dual deformation complex is due to Atiyah, Hitchin, and Singer
\cite{YM1954,AHS1978}.  The proofs below are elementary derivations of the
identities used here.  The analytic index formula and reconstruction theorems
are identified separately as external results.

Let $M$ be a smooth, oriented Riemannian manifold without boundary and let
$G$ be a compact Lie group.  Let $P\to M$ be a principal $G$-bundle and write
$\operatorname{ad}P=P\times_{\operatorname{Ad}}\mathfrak g$.  Choose an
$\operatorname{Ad}$-invariant positive definite inner product on
$\mathfrak g$.  If $G$ is realized as a compact matrix group in $U(m)$, a
positive multiple of $-\operatorname{Re}\operatorname{Tr}(XY)$ is an example
on its skew-Hermitian Lie algebra.  The numerical multiple matters when
comparing coupling constants, instanton charges, and short-distance
coefficients.  It does not change the differential identities.

The affine space of smooth connections has tangent space
$\Omega^1(M;\operatorname{ad}P)$.  In a local trivialization a connection is
a Lie-algebra-valued one-form $A$, with
\begin{equation}
 F_A=dA+\tfrac12[A\wedge A],\qquad
 d_A\alpha=d\alpha+[A\wedge\alpha].
 \label{eq:curvature-conventions}
\end{equation}
For homogeneous forms the bracket includes exterior multiplication and the
Lie bracket.  In coordinates this means
\[
 (F_A)_{\mu\nu}=\partial_\mu A_\nu-\partial_\nu A_\mu+[A_\mu,A_\nu],
 \qquad D_\mu X=\partial_\mu X+[A_\mu,X].
\]
The bracket is graded antisymmetric:
$[\alpha\wedge\beta]=-(-1)^{pq}[\beta\wedge\alpha]$ for degrees $p,q$.
In particular it is symmetric for two one-forms.  Keeping this sign is
essential in the second variation.

\begin{proposition}[Curvature is the failure of the complex identity]
For every adjoint-valued form $\alpha$,
\begin{equation}
 d_A^2\alpha=[F_A\wedge\alpha],\qquad d_AF_A=0.
 \label{eq:curvature-square}
\end{equation}
Consequently, if $F_A=0$, the sequence
$\Omega^0(\operatorname{ad}P)\xrightarrow{d_A}
\Omega^1(\operatorname{ad}P)\xrightarrow{d_A}\cdots$ is a complex.
\end{proposition}

\begin{proof}
For a degree-zero section $s$, direct expansion gives
\[
 (D_\mu D_\nu-D_\nu D_\mu)s
 =[\partial_\mu A_\nu-\partial_\nu A_\mu,s]
   +[A_\mu,[A_\nu,s]]-[A_\nu,[A_\mu,s]]
 =[F_{\mu\nu},s],
\]
where the last equality is the Jacobi identity.  Antisymmetrizing the same
calculation in the additional form indices gives the first identity for every
degree; differentiation of the scalar form coefficients contributes the
ordinary identity $d^2=0$.  Equivalently, the graded Leibniz rule and graded
Jacobi identity give
$d_A^2=[dA+\frac12[A\wedge A],\cdot]$ directly.
For the second identity, the cyclic sum of
$D_\lambda F_{\mu\nu}$ has cancelling second derivatives and cancelling
first-derivative brackets.  The remaining terms are
$[A_\lambda,[A_\mu,A_\nu]]$ and their cyclic permutations, whose sum is
zero by Jacobi.  This is precisely $d_AF_A=0$.  If $F_A=0$, the first
identity says that every consecutive composition is zero.
\end{proof}

One must not reverse the last implication without a hypothesis on the
representation.  A central curvature can act trivially in the adjoint
representation.  For instance an abelian connection may have nonzero
curvature although its induced adjoint differential is the ordinary exterior
derivative.  For a centerless Lie algebra the statement that
$[F_A,\phi]=0$ for every local adjoint section $\phi$ does force $F_A=0$.
This qualification is a small example of the semantic checks required when a
formal statement is used as a bridge between different mathematical settings.

\section{Gauge covariance and the variational equations}

Use the convention
$A^u=u^{-1}Au+u^{-1}du$ for a gauge transformation $u$.
Then $d_{A^u}(u^{-1}\alpha u)=u^{-1}(d_A\alpha)u$ and
$F_{A^u}=u^{-1}F_Au$.  These formulas can be checked in a local matrix
representation by differentiating $u^{-1}u=1$.  Since the chosen inner
product is invariant, the Yang--Mills functional
\begin{equation}
 S(A)=\frac{1}{2g^2}\int_M |F_A|^2\,d\operatorname{vol}
 \label{eq:yang-mills-action}
\end{equation}
is gauge invariant.  Here $g>0$ is fixed.  On a noncompact manifold all
variations used below have compact support, or sufficient decay to justify
integration by parts.  On a manifold with boundary a boundary condition must
be supplied; omitting that condition changes the Euler--Lagrange problem.

\begin{proposition}[First and second variations]
For $a,b\in\Omega^1(M;\operatorname{ad}P)$,
\begin{align}
 F_{A+ta}&=F_A+t\,d_Aa+\tfrac12t^2[a\wedge a],\label{eq:curvature-variation}\\
 DS_A(a)&=g^{-2}\langle d_A^*F_A,a\rangle_{L^2},\label{eq:ym-first}\\
 D^2S_A(a,b)&=g^{-2}\left(
 \langle d_Aa,d_Ab\rangle_{L^2}
 +\langle F_A,[a\wedge b]\rangle_{L^2}\right).
 \label{eq:ym-second}
\end{align}
Thus a smooth connection is stationary for all permitted variations if and
only if $d_A^*F_A=0$.
\end{proposition}

\begin{proof}
Substitute $A+ta$ into \eqref{eq:curvature-conventions}.  The linear terms
are $da+[A\wedge a]=d_Aa$ because the bracket of one-forms is symmetric.
The quadratic term is $\frac12t^2[a\wedge a]$.  Differentiating
$\frac12g^{-2}\langle F,F\rangle$ once gives
$g^{-2}\langle F_A,d_Aa\rangle$, which equals
$g^{-2}\langle d_A^*F_A,a\rangle$ by the definition of the formal
adjoint.  Mixed differentiation of $F_{A+sa+tb}$ gives
$\partial_s\partial_tF=[a\wedge b]$.  Applying the product rule to the
norm square proves the Hessian formula.  If the first variation vanishes
for every compactly supported $a$, the fundamental lemma of the calculus of
variations gives the stated equation; its converse follows from the formula.
\end{proof}

The equation $d_AF_A=0$ is an identity, while $d_A^*F_A=0$ is a field
equation.  Their different logical roles are important.  Imposing the first
does not restrict the set of connections, whereas imposing the second
selects the classical Yang--Mills critical points.  In a Lean development the
Bianchi identity should be proved from the definition of curvature, not
added as an unexplained assumption that could conceal an inconsistent
definition.

The infinitesimal gauge direction is $d_A\phi$, because
$\frac{d}{dt}|_{t=0} A^{\exp(t\phi)}=d_A\phi$.  At a critical point,
the Hessian annihilates these directions in its associated bilinear form.
One way to see this without a coordinate calculation is to differentiate
the gauge invariance identity twice along a gauge orbit and a transverse
variation.  The term involving acceleration of the orbit is paired with
$DS_A$, which is zero at a critical point.  Thus the degeneracy is forced by
symmetry, and the unmodified Hessian cannot be elliptic on all one-forms.

For a nonzero covector $\xi$, omit the harmless factor $i$ in the symbol
convention.  The principal symbol of $d_A$ is exterior multiplication by
$\xi$; that of its adjoint is contraction by $\xi^\sharp$, up to the
consistent sign for a first-order adjoint.  The principal symbol of the
nonnegative second-order operator $d_A^*d_A$ on one-forms is
\[
 a\longmapsto |\xi|^2a-\xi\,\iota_{\xi^\sharp}a.
\]
It vanishes on multiples of $\xi$.  Adding the gauge-fixing term
$d_Ad_A^*$ changes this to $|\xi|^2a$.  Curvature contributes only
lower-order terms to the Hessian.  Accordingly the gauge-fixed Hessian is
elliptic, although its lower-order part need not be positive at an arbitrary
critical point.  Ellipticity concerns invertibility of the principal symbol;
positivity is a separate property of the complete operator.

\section{Four dimensions and self-duality}

Suppose now that $\dim M=4$.  On two-forms the Hodge star satisfies
$*^2=1$, and hence
\[
 \Lambda^2=\Lambda^2_+\oplus\Lambda^2_-,\qquad
 F_A^\pm=\tfrac12(F_A\pm *F_A).
\]
The two summands are orthogonal.  Define
$Q(A)=\int_M\langle F_A\wedge F_A\rangle$, where the invariant
pairing is extended to forms.  Then
\begin{equation}
 \|F_A\|_2^2=\|F_A^+\|_2^2+\|F_A^-\|_2^2,
 \qquad Q(A)=\|F_A^+\|_2^2-\|F_A^-\|_2^2.
 \label{eq:energy-topology}
\end{equation}
For a closed manifold, $Q(A)$ is constant as the connection varies on a
fixed bundle.  Its normalization as an integral characteristic number depends
on the representation and the invariant form.

\begin{proposition}[Energy bound and classical minimizers]
With the conventions above,
$S(A)\geq |Q(A)|/(2g^2)$.  Equality holds exactly when one of
$F_A^+$ or $F_A^-$ vanishes.  Every such connection satisfies the
Yang--Mills equation.
\end{proposition}

\begin{proof}
For nonnegative $x,y$, $x+y\geq|x-y|$, with equality if and only if
$xy=0$.  Apply this to the squared norms in
\eqref{eq:energy-topology}.  A smooth form with zero $L^2$ norm is zero,
so equality is the self-dual or anti-self-dual condition.  If
$*F_A=\pm F_A$, then $d_A*F_A=\pm d_AF_A=0$ by Bianchi.  The
formula for the formal adjoint, $d_A^*=\pm *d_A*$ in the relevant
degree, therefore gives $d_A^*F_A=0$.
\end{proof}

For completeness, the connection independence of $Q$ follows directly from
the first variation.  Differentiation gives
$2\int_M\langle d_Aa\wedge F_A\rangle$.
Invariance of the pairing and Bianchi identify the integrand with
$d\langle a\wedge F_A\rangle$.  Stokes' theorem makes the integral
zero on a closed manifold.  The statement requires an additional boundary
analysis on noncompact spaces; finite action alone should not be silently
substituted for all the boundary hypotheses needed in a particular theorem.

This is a useful first encounter with a topological quantity controlling an
analytic variational problem.  It is nevertheless a statement about
classical connections.  It gives neither an interacting quantum Hilbert
space nor a positive lowest nonzero quantum energy.  In particular a
positive action of a nontrivial instanton is not the mass of a lowest
particle.  The action of a Euclidean field configuration and the spectrum
of a reconstructed Hamiltonian are different objects.

\section{The anti-self-dual deformation complex}

At an anti-self-dual connection, $F_A^+=0$.  Linearization of this equation
gives $d_A^+a=0$, where $d_A^+=\pi_+d_A$.  Quotienting by
infinitesimal gauge transformations leads to
\begin{equation}
 0\longrightarrow\Omega^0(\operatorname{ad}P)
 \xrightarrow{d_A}\Omega^1(\operatorname{ad}P)
 \xrightarrow{d_A^+}\Omega^2_+(\operatorname{ad}P)
 \longrightarrow0.
 \label{eq:asd-complex}
\end{equation}
The following proof gives both parts of the assertion that this is an
elliptic complex.  Having the correct ranks is not by itself sufficient:
one needs the composition identity and exactness of the symbol at every
nonzero cotangent vector.

\begin{theorem}[Exactness of the anti-self-dual symbol sequence]
If $F_A^+=0$, then $d_A^+d_A=0$.  For every nonzero covector $\xi$,
the symbol sequence of \eqref{eq:asd-complex} is exact.
\end{theorem}

\begin{proof}
For an adjoint section $\phi$, the curvature identity gives
$d_A^+d_A\phi=\pi_+[F_A,\phi]=[F_A^+,\phi]=0$.
For symbol exactness the Lie algebra factor is passive, so it suffices to
work with ordinary real forms.  Rescale and choose an oriented orthonormal
coframe $e^1,e^2,e^3,e^4$ with $\xi=|\xi|e^1$.  The first symbol sends
$c$ to $c\xi$ and is injective.  If
$a=\sum_j a_je^j$, the second sends it to
\[
 \frac{|\xi|}{2}\sum_{j=2}^4a_j
 \left(e^1\wedge e^j+*(e^1\wedge e^j)\right).
\]
The three parenthesized forms are mutually orthogonal and nonzero, and
they span the three-dimensional self-dual space.  The map is therefore
surjective and its kernel consists precisely of multiples of $e^1$.
This kernel is the image of the first symbol.  Tensoring the exact
finite-dimensional sequence with $\mathfrak g$ preserves exactness.
\end{proof}

The proof displays a direct algebraic connection to varieties of complexes:
each nonzero cotangent vector determines a pair of linear maps satisfying a
zero-product equation and a particular exact rank condition.  The geometric
content is the existence of a smooth family of these maps, with specified
homogeneity in $\xi$, over the punctured cotangent bundle.  A single
algebraic complex cannot replace that family.  Its global symbol class is
the input to index theory.

Let $H_A^0,H_A^1,H_A^2$ denote the cohomology of
\eqref{eq:asd-complex}.  On a closed manifold elliptic theory makes these
spaces finite dimensional and gives harmonic representatives after choosing
Sobolev completions.  Their meanings are distinct: $H_A^0$ contains
infinitesimal stabilizers, $H_A^1$ infinitesimal deformations modulo gauge,
and $H_A^2$ obstructions.  The operator
\[
 L_A=d_A^*\oplus d_A^+:
 \Omega^1(\operatorname{ad}P)\longrightarrow
 \Omega^0(\operatorname{ad}P)\oplus\Omega^2_+(\operatorname{ad}P)
\]
has index $\dim H_A^1-\dim H_A^0-\dim H_A^2$.
At an irreducible, unobstructed connection this is the local dimension of
the moduli space.  Without those conditions the same integer is a virtual
dimension, not an assertion that the quotient is a manifold.

\begin{theorem}[Atiyah--Hitchin--Singer index formula, external]
For an $SU(2)$ bundle on a closed, connected, oriented Riemannian four-manifold,
write $k=\langle c_2(P\times_{SU(2)}\mathbb C^2),[M]\rangle$ with
the convention in which anti-self-dual instantons have positive $k$.
Then
\begin{equation}
 \operatorname{ind}L_A=8k-3(1-b_1(M)+b_2^+(M)).
 \label{eq:asd-index}
\end{equation}
\end{theorem}

This external theorem is an application of the analytic index theorem to
the symbol just computed \cite{AHS1978}.  The proof of the symbol
calculation above is not a proof of the index theorem.  On $S^4$ the
formula becomes $8k-3$ for the unframed irreducible problem.  Framing at
a point removes the three-dimensional constant gauge quotient in the usual
regular situation and gives $8k$.  Comparing these numbers without
recording framing conventions creates a spurious disagreement.

\section{A finite-dimensional algebraic model: monads and ADHM}

There is a stronger bridge than a resemblance between equations.  The
Atiyah--Drinfeld--Hitchin--Manin construction relates appropriate
finite-dimensional algebraic data to framed instantons
\cite{ADHM1978}.  We use only its elementary complex identity here, and
state the correspondence as an external theorem rather than re-proving its
analytic and algebro-geometric foundations.

Let $V,W$ be finite-dimensional complex vector spaces.  Choose
$B_1,B_2\in\operatorname{End}(V)$,
$I\in\operatorname{Hom}(W,V)$, and
$J\in\operatorname{Hom}(V,W)$.  Set
\[
 \mu_{\mathbb C}=[B_1,B_2]+IJ.
\]
For $z=(z_1,z_2)\in\mathbb C^2$, define
\[
 \alpha_z(v)=((B_1-z_1)v,(B_2-z_2)v,Jv),
 \quad
 \beta_z(v_1,v_2,w)=-(B_2-z_2)v_1+(B_1-z_1)v_2+Iw.
\]
Their types are $V\to V\oplus V\oplus W\to V$.  Direct
multiplication gives
\begin{equation}
 \beta_z\alpha_z=[B_1,B_2]+IJ=\mu_{\mathbb C}.
 \label{eq:adhm-composition}
\end{equation}
Thus the complex moment-map equation is exactly a zero-composition
condition, valid simultaneously for all $z$.

If $\alpha_z$ is injective and $\beta_z$ surjective for every relevant
$z$, the middle cohomology has dimension $\dim W$.  Indeed
$\dim\ker\beta_z=2\dim V+\dim W-\dim V$, and quotienting
by the $\dim V$-dimensional image of $\alpha_z$ leaves $\dim W$.
These constant-rank conditions give a vector bundle locally.  They are
essential: the equation $\mu_{\mathbb C}=0$ alone permits rank drops
and singular objects.

For unitary instantons the real moment-map equation also enters:
\[
 [B_1,B_1^*]+[B_2,B_2^*]+II^*-J^*J=0.
\]
Together with the appropriate stability and regularity conditions, the
data modulo the unitary action on $V$ yield the ADHM correspondence.
The complex equation explains why varieties of complexes are relevant;
the real equation, positivity, and rank conditions explain why arbitrary
points of a variety of complexes are insufficient.  In particular this
construction describes classical instantons.  Its finite-dimensional
algebra does not prove a quantum mass gap or an algorithm for elliptic-curve
discrete logarithms.

The monad identity is also a good candidate for machine checking: matrix
multiplication, commutators, scalar identities, and dimensions have clear
finite specifications.  Formalizing the full ADHM correspondence would
require extensive differential geometry and analysis beyond that elementary
identity.  Reporting the identity as ``the ADHM theorem verified'' would
discard almost all of the theorem's content.

\section{Why an index cannot establish a quantitative gap}

An index records a difference of dimensions.  A spectral gap is an
inequality with a positive numerical constant.  The distinction can be
proved by an example that is already a complex of finite-dimensional
vector spaces.

\begin{proposition}[Exactness does not control the smallest singular value]
For every $\varepsilon>0$, the complex
$0\to\mathbb R\xrightarrow{\varepsilon}\mathbb R\to0$ is exact,
but the smallest eigenvalue of either associated nonzero Laplacian is
$\varepsilon^2$.  There is no positive lower bound valid for all
$\varepsilon>0$.
\end{proposition}

\begin{proof}
Multiplication by a nonzero real number is a bijection, so both cohomology
groups vanish.  Relative to the standard inner products, the adjoint is
also multiplication by $\varepsilon$, and $d^*d=dd^*=\varepsilon^2$.
For any proposed lower bound $c>0$, choosing
$0<\varepsilon<\sqrt c$ contradicts the bound.
\end{proof}

The same phenomenon appears geometrically.  On a flat torus of side length
$L$, the scalar Laplacian has eigenvalues
$(2\pi/L)^2|n|^2$ for $n\in\mathbb Z^d$.  The first positive one is
$4\pi^2/L^2$, which tends to zero as $L\to\infty$.  Each fixed torus
has a compact elliptic problem and a discrete spectrum, yet no
volume-independent positive gap follows.  The quantum Yang--Mills problem
contains a much more difficult version of precisely this issue: a bound
must survive both the ultraviolet limit and unbounded physical volume.

This gives a disciplined place for orders and homological algebra.  They
can organize integral models, modules, deformation complexes, characteristic
classes, and exact relations among observables.  To conclude a mass gap
one still needs coercive estimates or spectral information with constants
that remain positive in the required limits.  The absence of such an
estimate is a mathematical missing implication, not a matter of choosing
more persuasive terminology.

\chapter{From finite gauge lattices to a quantum spectral gap}
\label{chap:quantum-gap}

\section{Three different uses of the word lattice}

A lattice over an order is a module with a finiteness and torsion-freeness
condition over the ground ring.  A Euclidean lattice is a discrete subgroup
of a real inner-product space, with a metric that makes shortest-vector and
decoding questions meaningful.  A lattice in a lattice gauge theory is a
discrete approximation to spacetime, equipped with vertices, links, and
plaquettes.  The first two notions can be related by embeddings of number
fields and choices of positive forms.  The third need not carry any order
module structure at all.  Identifying all three because they share a word
would erase their defining mathematical data.

There are productive interactions nevertheless.  A gauge theory on a finite
cell complex has incidence maps and chain groups.  A choice of coefficient
ring can make those chain groups modules over an order.  Discrete Hodge
operators add positive inner products.  Fourier analysis on finite abelian
groups, central in quantum algorithms, also diagonalizes many free lattice
operators.  Each such connection is a construction that can be specified
and proved.  None entails that the computational hardness of a Euclidean
lattice problem controls the mass spectrum of an interacting gauge theory.

\section{A finite Wilson model exists by elementary integration}

Let a finite oriented cell complex have a set $E$ of positively oriented
edges and a set $\mathcal P$ of oriented plaquettes.  Assign $U_e\in G$ to
each $e\in E$, and use $U_{\bar e}=U_e^{-1}$ for the reverse orientation.
The ordered product around a plaquette is $U_p$.  For a fixed unitary
representation $G\subset U(m)$ and $\beta\geq0$, define
\begin{equation}
 E_W(U)=\sum_{p\in\mathcal P}
 \left(1-\frac1m\operatorname{Re}\operatorname{Tr}U_p\right),\quad
 Z_\beta=\int_{G^E}e^{-\beta E_W(U)}\prod_{e\in E}dh(U_e),
 \label{eq:wilson-finite}
\end{equation}
where $h$ is normalized Haar measure.  Wilson's lattice construction
preserves gauge invariance before taking any limit \cite{Wilson1974}.
The following modest result explains exactly how much existence is obtained
from compactness alone.

\begin{proposition}[Finite-volume normalization and gauge invariance]
The quantity $Z_\beta$ satisfies
$e^{-2\beta|\mathcal P|}\leq Z_\beta\leq1$, and
$d\mu_\beta=Z_\beta^{-1}e^{-\beta E_W}\prod_e dh(U_e)$ is a
probability measure.  It is invariant under the vertex gauge action
$U_{xy}\mapsto t_xU_{xy}t_y^{-1}$.
\end{proposition}

\begin{proof}
Every eigenvalue of a unitary matrix has modulus one, so
$-m\leq\operatorname{Re}\operatorname{Tr}U_p\leq m$.
Thus $0\leq E_W\leq2|\mathcal P|$ and
$e^{-2\beta|\mathcal P|}\leq e^{-\beta E_W}\leq1$.
Integration against a probability measure gives the bounds.  In particular
the normalization is finite and strictly positive.
Under a vertex transformation the successive $t_y^{-1}t_y$ factors in a
plaquette cancel, leaving conjugation by the transformation at its initial
vertex.  Trace is invariant under conjugation, so the action is unchanged.
Haar measure on a compact group is invariant under both left and right
translations.  Applying this separately to each edge preserves the product
measure and proves the last assertion.
\end{proof}

The proof works for every finite complex and every compact matrix group.
It contains no control as the number of plaquettes tends to infinity or as
the lattice spacing tends to zero.  The lower bound on $Z_\beta$ itself
degenerates exponentially with $|\mathcal P|$.  Dividing by $Z_\beta$
produces a probability measure, but does not remove the need for uniform
bounds on the observables whose limits one wants to take.

An observable such as a product of finitely many matrix entries is bounded
at fixed lattice spacing.  A renormalized curvature observable typically
contains negative powers of the spacing and cutoff-dependent mixing with
other observables.  Its bound can diverge as the spacing tends to zero.
Consequently, compactness of $G$ is compatible with difficult and highly
singular continuum limits.  The relevant estimates must concern the
renormalized observables and their correlations, not merely the unscaled
link variables.

\section{Reflection positivity is additional structure}

Let $\vartheta$ reflect Euclidean time and let $\mathcal A_+$ denote a
space of observables supported at positive times.  Write
$\Theta F=\overline{F\circ\vartheta}$, with any required adjoint on
field labels.  Reflection positivity asserts
\begin{equation}
 \mathbb E[(\Theta F)F]\geq0\qquad(F\in\mathcal A_+).
 \label{eq:reflection-positive}
\end{equation}
It is not a consequence of the ordinary positivity of a probability measure.
For example, take a symmetric random variable $Z$ with
$0<\mathbb E[Z^2]<\infty$ and set $X_+=Z$, $X_-=-Z$.
The joint law is a positive probability measure and is invariant under
exchanging the two time labels.  Nevertheless, for $F=X_+$,
$\mathbb E[(\Theta F)F]=\mathbb E[X_-X_+]=-\mathbb E[Z^2]<0$.
Both positivity and reflection symmetry can hold while reflection positivity
fails.

The Wilson action has the special structure needed for appropriate lattice
reflection-positivity and transfer-matrix theorems
\cite{OSSeiler1978,Luscher1977}.  A useful finite-dimensional ingredient
can be proved directly.

\begin{proposition}[Positive-type Wilson link weight]
Let $\rho:G\to U(m)$ be a unitary representation, $\chi(g)=\operatorname{Tr}\rho(g)$,
and $c\geq0$.  The function $w(g)=\exp(c\operatorname{Re}\chi(g))$
is of positive type: for all $g_1,\ldots,g_n\in G$ and
$z_1,\ldots,z_n\in\mathbb C$,
\[
 \sum_{i,j}\overline{z_i}z_jw(g_i^{-1}g_j)\geq0.
\]
\end{proposition}

\begin{proof}
The character $\chi$ is of positive type because
\[
 \sum_{i,j}\overline{z_i}z_j\chi(g_i^{-1}g_j)
 =\operatorname{Tr}\left(
  \Bigl(\sum_i z_i\rho(g_i)\Bigr)^*
  \Bigl(\sum_j z_j\rho(g_j)\Bigr)\right)\geq0.
\]
The conjugate character is the character of the conjugate unitary
representation.  Products of characters are characters of tensor products.
Expanding
\[
 w(g)=\sum_{a,b\geq0}\frac{(c/2)^{a+b}}{a!b!}
                  \chi(g)^a\overline{\chi(g)}^{b}
\]
therefore expresses $w$ as a limit of nonnegative linear combinations of
positive-type functions.  Since $|\chi(g)|\leq m$, the double series
converges uniformly.  For any fixed finite set of points and coefficients,
the associated nonnegative quadratic forms converge to the displayed
quadratic form, which remains nonnegative.
\end{proof}

This lemma proves a kernel positivity statement.  The full lattice theorem
also needs a factorization compatible with the reflection plane, the
orientation of links, and the integration of gauge variables.  Those
additional steps are not omitted by calling the lemma a proof of the full
Osterwalder--Schrader reconstruction theorem.  The separation of claims is
particularly useful in a formal development: a locally verified kernel
identity can be reused without overstating its global consequence.

\section{The Hilbert-space quotient behind reconstruction}

Reflection positivity supplies the semidefinite sesquilinear form
$(F,G)_{\mathrm{OS}}=\mathbb E[(\Theta F)G]$.
Assume that the correlation functional is defined on a complex vector
space $\mathcal V_+$ of finite positive-time field words, that this form
is Hermitian, and that it is nonnegative on the diagonal.  The following
quotient step is purely algebraic.

\begin{proposition}[Null vectors and positive quotient]
For a positive semidefinite Hermitian form $B$ on a complex vector space,
the set $N=\{v:B(v,v)=0\}$ is a linear subspace and
$B(v,w)=0$ for every $v\in N$ and $w$.  The induced form on
$\mathcal V_+/N$ is positive definite.
\end{proposition}

\begin{proof}
If $B(w,w)>0$, positivity of
$B(v+\lambda w,v+\lambda w)$ for every $\lambda\in\mathbb C$
and minimization of this real quadratic polynomial gives
$|B(v,w)|^2\leq B(v,v)B(w,w)$.  If $B(w,w)=0$, apply the same
argument to $w+\varepsilon u$ with a positive-norm $u$ and pass to the
limit, or examine the linear term of the polynomial directly.  Thus
Cauchy--Schwarz holds in all cases.  If $v\in N$ it forces
$B(v,w)=0$.  This orthogonality implies that sums and scalar multiples of
null vectors are null, and that changing a representative by a null vector
does not change any inner product.  Finally an equivalence class with zero
norm has a representative in $N$ and is therefore the zero class.
\end{proof}

Completing this quotient gives a Hilbert space.  Euclidean time translation,
under the additional continuity and boundedness estimates required in
reconstruction, induces a strongly continuous self-adjoint contraction
semigroup.  Its generator is a nonnegative self-adjoint operator $H$ and
the semigroup is $e^{-tH}$.  A complete reconstruction must also supply
field domains, distributional regularity, locality, analytic continuation, and the Poincare action.  The original Osterwalder--Schrader papers state
precise regularity and growth hypotheses for these conclusions
\cite{OS1973,OS1975}.  The quotient lemma alone supplies none of them.

One should track test-function spaces carefully.  A correlation kernel
defined only where $x_i\ne x_j$ is not automatically a tempered
distribution on all configurations.  Terms supported on the coincidence
diagonals are invisible on separated tests.  A simple illustration is
$T$ and $T+c\delta_0$ on $\mathbb R^d$: they agree on every test
function supported away from the origin, but disagree on a test with
$f(0)\ne0$.  Even all derivatives of a test vanishing on a diagonal
make every derivative of the corresponding delta distribution invisible.
This does not invalidate a separated-point construction.  It specifies the
extra work necessary to extend it, establish positivity of an extension,
and prove uniqueness under a chosen normalization.

\section{What the mass-gap statement actually requires}

In the Hamiltonian formulation, fix a complex Hilbert space $\mathcal H$,
a normalized vacuum $\Omega$, and a nonnegative self-adjoint Hamiltonian
$H$ with $\ker H=\mathbb C\Omega$.  Let $P_0$ be orthogonal
projection onto the vacuum.  A positive gap means that there exists
$\Delta>0$ with
\begin{equation}
 \operatorname{spec}(H)\subset\{0\}\cup[\Delta,\infty).
 \label{eq:mass-gap}
\end{equation}
For the Yang--Mills Millennium problem this spectral assertion accompanies
a nontrivial four-dimensional quantum field theory with the required
Yang--Mills local-field and short-distance structure.  It is not an
independent name for any positive matrix eigenvalue
\cite{JaffeWitten2006}.  Vacuum uniqueness, positive energy, existence of a
nonzero local observable, and a mass gap are four different conditions.

\begin{theorem}[Spectral gap and Euclidean-time decay]
Under the hypotheses just stated, the following are equivalent for a fixed
$\Delta>0$:
\begin{enumerate}
 \item \eqref{eq:mass-gap} holds.
 \item $\|e^{-tH}(1-P_0)\|\leq e^{-t\Delta}$ for every $t\geq0$.
 \item For every $v\perp\Omega$ and every $t\geq0$,
 $0\leq\langle v,e^{-tH}v\rangle\leq e^{-t\Delta}\|v\|^2$.
\end{enumerate}
\end{theorem}

\begin{proof}
Let $E_H$ be the spectral measure of $H$.  Under the first condition,
the restriction to $\Omega^\perp$ has spectral support in
$[\Delta,\infty)$.  Functional calculus bounds
$e^{-t\lambda}$ by $e^{-t\Delta}$ on that support, proving the
second condition.  Positivity of $e^{-tH}$ and Cauchy--Schwarz then give
the third.

Conversely suppose the third condition and suppose the spectral projection
$E_H((0,\Delta))$ is nonzero.  Since $(0,\Delta)$ is a countable
union of closed intervals $[a,b]$ with $0<a\leq b<\Delta$, one such
projection is nonzero.  Choose a nonzero $v$ in its range.  It is
orthogonal to the zero-energy eigenspace, and
\[
 \langle v,e^{-tH}v\rangle
 =\int_{[a,b]}e^{-t\lambda}\,d\langle v,E_H(\lambda)v\rangle
 \geq e^{-tb}\|v\|^2.
\]
For every $t>0$, $e^{-tb}>e^{-t\Delta}$, contradicting the third
condition.  Thus the projection is zero.  The spectral theorem then gives
the first condition.  The case $t=0$ is the equality bound and causes no
exception.
\end{proof}

For a bounded observable $A$, set
$a=(A-\langle\Omega,A\Omega\rangle)\Omega$.
Then $a\perp\Omega$, and the associated Euclidean-time connected
two-point function is
\[
 C_A(t)=\langle a,e^{-tH}a\rangle.
\]
The theorem bounds it by $e^{-t\Delta}\|a\|^2$.
For unbounded smeared fields, the same statement applies once the vacuum
is in the relevant operator domain and $a$ is a Hilbert-space vector.
Writing products of unsmeared fields at the same point would evade this
domain issue, not solve it.  Spatial clustering in a relativistic theory
requires the locality and spectrum properties that relate time and space;
the elementary theorem here concerns Euclidean time only.

A useful converse needs control on a total collection of vectors.  If the
third inequality holds with one common $\Delta$ on a linear dense
subspace of $\Omega^\perp$, continuity of the bounded operator
$e^{-tH}$ extends it to every vector, and the theorem proves the gap.
Control of a single observable cannot do this.  For example take
$\mathcal H=\mathbb C\Omega\oplus\mathbb C a\oplus L^2(0,1)$,
let $H\Omega=0$, $Ha=a$, and let $H$ act by multiplication by
$x$ on $L^2(0,1)$.  The $a$-correlation is exactly
$e^{-t}\|a\|^2$, but the third summand has spectrum down to zero.
The complete Hamiltonian has no positive mass gap.

\section{A rigorous countermodel to ultraviolet-to-infrared inference}

Short-distance asymptotics probe high spectral energies, whereas a gap
concerns low energies.  This distinction can be shown with two explicit
positive spectral measures.  Define for $t>0$
\begin{align}
 C_{\mathrm{g}}(t)&=\int_1^\infty
       \frac{E^7}{\log^2(e+E)}e^{-tE}\,dE,\\
 C_{\mathrm{ng}}(t)&=C_{\mathrm{g}}(t)+\int_0^1e^{-tE}\,dE.
 \label{eq:spectral-countermodel}
\end{align}
The first measure has support in $[1,\infty)$; the second has support
arbitrarily close to zero.  Both Laplace transforms are finite for every
$t>0$ and have the same leading short-distance behavior.

\begin{proposition}[Identical leading ultraviolet behavior, different gaps]
For either function in \eqref{eq:spectral-countermodel},
\[
 \lim_{t\downarrow0}t^8\log^2(1/t)C(t)=7!.
\]
The corresponding positive spectral measures have, respectively, a positive
lower support bound and no positive lower support bound.
\end{proposition}

\begin{proof}
Write $L=\log(1/t)$ and substitute $u=tE$ in the first integral.
After multiplication by $t^8L^2$, the result is
\[
 \int_t^\infty u^7e^{-u}
       \left(\frac{L}{\log(e+u/t)}\right)^2du.
\]
For $u\geq\sqrt t$ and $t<1$,
$\log(e+u/t)\geq L/2$, so the integrand is at most
$4u^7e^{-u}$.  For fixed $u>0$ the parenthesized ratio tends to one.
Dominated convergence on the region $u\geq\sqrt t$, extended by
zero elsewhere, therefore gives $\int_0^\infty u^7e^{-u}du=7!$.
The contribution from $t\leq u\leq\sqrt t$ is at most
$L^2\int_0^{\sqrt t}u^7du=L^2t^4/8$, which tends to zero.
This proves the first limit.  The additional integral is bounded by one,
and $t^8\log^2(1/t)\to0$, so it does not change the limit.
The support assertions follow immediately from the strictly positive
densities on their respective intervals.
\end{proof}

These examples are positive spectral models, not constructions of local
four-dimensional gauge theories.  Their role is logical: a proposed
inference using only positivity and this ultraviolet asymptotic cannot
force a gap, because the two examples satisfy those inputs and have
different infrared support.  Additional Yang--Mills-specific estimates may
well force a gap.  They must be stated and proved rather than inferred from
the leading asymptotic alone.

\section{Uniform estimates that would survive a limit}

Here is a precise version of one sufficient missing step.  Work on one
Hilbert space so that convergence has an unambiguous meaning.  Let $H_n$
and $H$ be nonnegative self-adjoint operators, and let $P_n,P$ be their
zero-energy spectral projections.  Assume
$P_n\to P$ strongly and
$e^{-tH_n}\to e^{-tH}$ strongly for each $t\geq0$.

\begin{proposition}[Passage of a uniform gap bound]
If for one $\Delta>0$ and all $n,t\geq0$,
\[
 \|e^{-tH_n}(1-P_n)\|\leq e^{-t\Delta},
\]
then $\|e^{-tH}(1-P)\|\leq e^{-t\Delta}$ for all $t\geq0$.
\end{proposition}

\begin{proof}
For fixed $v$ and $t$, write
\begin{align*}
 e^{-tH_n}(1-P_n)v-e^{-tH}(1-P)v
 &=e^{-tH_n}(P-P_n)v\\
 &\quad+(e^{-tH_n}-e^{-tH})(1-P)v.
\end{align*}
The first term tends to zero since the semigroups are contractions and
$P_n\to P$ strongly.  The second tends to zero by the assumed strong
convergence.  Taking norms in the uniform bound and passing to the limit
gives $\|e^{-tH}(1-P)v\|\leq e^{-t\Delta}\|v\|$.
Taking the supremum over unit vectors proves the claim.
\end{proof}

For actual lattice-to-continuum constructions the Hilbert spaces may vary.
One must first supply compatible embeddings or formulate convergence in
terms of a common observable algebra and its states.  The proposition is
a conditional theorem with clear hypotheses; it does not assert that
Yang--Mills satisfies those hypotheses.  It is useful because it localizes
the hard work in the uniform bound and the comparison maps, rather than
hiding it inside the word ``limit.''

Physical units are equally important.  If a transfer matrix $T_a$ is
normalized to have largest eigenvalue one and corresponds to time step
$a$, its Hamiltonian is $H_a=-a^{-1}\log T_a$ on the relevant
spectral subspace.  A bound for a second eigenvalue must therefore be
interpreted through $-a^{-1}\log\lambda_2(a)$.  Merely proving
$\lambda_2(a)<1$ separately for each $a$ does not imply a common
positive physical gap.  For example
$\lambda_2(a)=e^{-a^2}$ is strictly less than one but gives physical
gap $a\to0$.  This elementary counterexample rules out a frequent
confusion between finite-cutoff positivity and continuum spectral control.

\section{Renormalization and what a running coupling can prove}

The passage from the finite Wilson law to continuum fields requires more
than refinement of a mesh.  The bare coupling and the observable
normalizations must change with the cutoff.  The four-dimensional
renormalization-group program of Ba{\l}aban constructs effective actions
and controls transformations between scales; the 1987 paper cited here
explicitly works in its small-field approximation
\cite{Balaban1987}.  The limitation matters: a small-field expansion is
not automatically control of the full probability measure, and estimates
for all fields require additional arguments.

The original perturbative discovery of asymptotic freedom determines a
negative leading beta function for the appropriate nonabelian theories
\cite{GrossWilczek1973}.  An exact scalar differential equation captures
the leading behavior without claiming to be the full quantum theory.
Let $b>0$ and suppose, on an interval of scales $\mu>0$, that
\begin{equation}
 \mu\frac{dg}{d\mu}=-b g^3,\qquad g(\mu)>0.
 \label{eq:one-loop-ode}
\end{equation}
Then
\begin{equation}
 \frac1{g(\mu)^2}
 =\frac1{g(\mu_0)^2}+2b\log\frac\mu{\mu_0},
 \qquad
 g(\mu)^2=\frac1{2b\log(\mu/\Lambda)},
 \label{eq:running-coupling-toy}
\end{equation}
where
$\Lambda=\mu_0\exp(-1/(2b g(\mu_0)^2))$ and the displayed positive
solution is considered only for $\mu>\Lambda$.  The proof is simply
$\frac{d}{d\log\mu}(g^{-2})=-2g^{-3}\frac{dg}{d\log\mu}=2b$,
followed by integration.  In particular $g(\mu)\to0$ at large
$\mu$.

The parameter $\Lambda$ illustrates dimensional transmutation: an
initial dimensionless coupling together with a reference scale determines
a dimensional scale.  Equation \eqref{eq:running-coupling-toy} does not
prove a particle of mass $\Lambda$, a mass gap proportional to
$\Lambda$, or a controlled continuation through $\mu=\Lambda$.
Its denominator vanishes there because the leading differential model
leaves its useful weak-coupling range.  Deducing the infrared spectrum
from that singularity would replace a missing theorem with an artifact of
an approximation.

The same distinction appears in the control of irrelevant interactions.
For a local operator of engineering degree $d$ in four dimensions, the
classical rescaling factor of its coupling under a length enlargement by
$b_\mathrm{s}>1$ is $b_\mathrm{s}^{4-d}$.  The exponent explains the
words relevant, marginal, and irrelevant.  It does not bound the nonlinear
errors or the infinite collection of generated operators.  A simple
inequality isolates one of the required analytic mechanisms.

\begin{proposition}[A contractive recurrence with controlled errors]
Let $0\leq\rho<1$, $a_n\geq0$, $\varepsilon_n\geq0$, and
$a_{n+1}\leq\rho a_n+\varepsilon_n$.  Then
\[
 a_n\leq\rho^n a_0+
       \sum_{j=0}^{n-1}\rho^{n-1-j}\varepsilon_j.
\]
If $\varepsilon_j\leq\varepsilon$ for all $j$, this is at most
$\rho^n a_0+\varepsilon/(1-\rho)$.  If instead
$\varepsilon_j\to0$, then $a_n\to0$.
\end{proposition}

\begin{proof}
The finite-sum bound follows by induction.  Multiplying the bound for
$a_n$ by $\rho$ and adding $\varepsilon_n$ gives the bound for
$a_{n+1}$.  The uniform-error assertion follows from summing a geometric
series.  For the final assertion fix $\eta>0$ and choose $J$ with
$\varepsilon_j<(1-\rho)\eta/2$ for $j\geq J$.  The sum over these
indices is at most $\eta/2$.  The finitely many earlier terms, together
with $\rho^n a_0$, tend to zero.  For sufficiently large $n$ they
contribute less than $\eta/2$.  Hence $a_n<\eta$ eventually.
\end{proof}

A renormalization proof must identify the norms represented by $a_n$,
show that the nonlinear transformation satisfies the recurrence, and bound
the errors with the right uniformity in volume and insertion degree.
The elementary lemma makes none of those analytic inputs disappear.
What it supplies is a reusable logical step once those inputs exist.
This is also a sensible formalization boundary: verify the recurrence
lemma independently, then connect it to the actual effective-action
estimates through a theorem whose hypotheses remain visible.

\section{Pure phases, clustering, and the order of quantifiers}

Suppose a translation-invariant state is a mixture
$\omega=\int_X\omega_\xi\,d\nu(\xi)$ of component states.
Even if every component clusters, the mixture need not.  For a self-adjoint
observable $A$, assume enough boundedness to pass the large-separation
limit through the integral.  Component clustering gives
\[
 \lim_{|x|\to\infty}\omega(A\,\tau_x(A))
 =\int_X\omega_\xi(A)^2\,d\nu(\xi).
\]
Subtracting $\omega(A)^2$ leaves the variance of
$\xi\mapsto\omega_\xi(A)$, which can be strictly positive.
The argument is an elementary proof that a phase decomposition and
clustering of each phase do not imply clustering of the original mixture.

Likewise, the statement that every diverging volume sequence has a
convergent subsequence does not assert that all subsequential limits agree.
The scalar sequence $(-1)^n$ already distinguishes sequential compactness
from convergence.  In a field theory the extracted object must include all
correlation functions coherently.  Choosing a different subsequence for
each observable is weaker than finding one subsequence working for all of
them.  Conversely, finding one common subsequence remains weaker than
proving independence of the subsequence.  These distinctions are precisely
the ones needed to read the recent formal construction correctly.

There is also a separate distinction between uniqueness of a vacuum inside
one representation and uniqueness of the theory selected by a continuum
limit.  Two different pure phases can each have a unique vacuum in their
own Hilbert spaces.  Neither uniqueness statement identifies the phases
with one another.  A formal verification plan should represent these
questions with different predicates so that a theorem about one cannot be
mistakenly reported as a theorem about the other.

\chapter{The 2026 OpenAI Yang--Mills construction: statement and verification}
\label{chap:openai-ym}

\section{A dated research result, not the older amplitude result}

This chapter records the public evidence available on 13 November 2026.
OpenAI's announcement dated 26 October 2026 includes a claimed construction
of four-dimensional pure Yang--Mills theory, as part of broader announced
work on gauge theories with fermions.  The announcement identifies a
natural-language Yang--Mills existence proof dated 9 September and a
completed formalization dated 16 October.  It describes an independent
group organized around David Gross, with Michael Douglas, Arthur Jaffe,
Igor Klebanov, and Nikita Nekrasov, to examine and explain the work
\cite{OpenAICommunity2026}.  These are dates and claims of the announcing
organization; the subsequent released submission and the independent
assessment provide the mathematical objects examined here.

The independent group's assessment, dated 10 November 2026, addresses the
submission of 28 October.  Its stated judgment is that the submitted
construction yields a Wightman quantum field theory from a four-dimensional
Wilson lattice scaling limit for the specified simple gauge groups.  It
explicitly withholds a mass-gap claim and uniqueness of the infinite-volume
limit.  The assessors also say they do not yet fully understand the proof
and distinguish their review from a Clay Mathematics Institute decision
\cite{YMReview2026}.  It is therefore inaccurate both to dismiss the
announcement as merely a perturbative amplitude calculation and to present
it as a completed solution of the whole Millennium problem.

The earlier paper on single-minus gluon amplitudes is a different result.
Its arXiv submission was posted on 12 February 2026, with a second version
on 9 March.  It concerns special half-collinear kinematics in Klein space
or complexified momenta, rather than a four-dimensional nonperturbative
vacuum construction \cite{Guevara2026}.  Its AI-assisted derivation may be
relevant to research methodology, but it does not supply the infrared
estimate required in \eqref{eq:mass-gap}.  Conflating the two projects would
misstate both the scope and the chronology of the recent work.

\section{Immutable source identities and review coverage}

The original public source was inspected at commit
\texttt{3456e5e75cf8d4cd558955a5c7638fba2d68f464} of
\texttt{yang-mills-group/OpenAI-YM-20261028}.  The group's review material
was inspected at commit
\texttt{55ba9fc36293b3f4d52c6b1f813169a0be76ee98} of
\texttt{yang-mills-group/YM}.  The accompanying evidence manifest records
URLs and hashes of the files used in this review.  The 35-page assessment
was read in full, together with the headline theorem, its public
specifications, the stronger construction theorem, the combined acceptance
specification and proof, and selected explanatory verification modules.
This is a review of those materials, not a claim to have read every proof
in the approximately million-line upstream development.

The distinction matters operationally.  A source snapshot, a successful
local compilation, a kernel replay, and an independent semantic review are
different deliverables.  The present project can pin and analyze an
external formal result without pretending that a historical upstream build
was performed again here.  Reproduction instructions must identify the
toolchain, dependencies, target declarations, and exact evidence actually
produced by a new run.  A link to an upstream pass is valuable evidence of
that upstream check; it is not a substitute for a fresh local pass when a
fresh pass is being claimed.

\section{The headline theorem and its quantifiers}

At the frozen revision, the file \texttt{Solution.lean} proves the
declaration \texttt{YangMills.jointExistence}.  Its target is defined in
\texttt{ChallengeSpec.lean}.  In ordinary mathematical notation, its
quantifier pattern is
\begin{equation}
 \forall G\;\bigl(\operatorname{Simple}(G)\Longrightarrow
 \exists C\,[\operatorname{Reg}(C)\ \wedge\
 \forall p\in\operatorname{Inc}(\mathbb N)
 \ \exists q\in\operatorname{Inc}(\mathbb N)
 \ \exists T\;\operatorname{Properties}(C,T,p\circ q)]\bigr).
 \label{eq:openai-quantifiers}
\end{equation}
Here the gauge-group type is a closed matrix subgroup of a unitary group,
and simplicity concerns its nonabelian real Lie algebra.  The object $C$
is one Wilson family, $p$ requests an increasing sequence of physical
volume indices, $q$ refines it, and $T$ is a continuum hierarchy associated
to the same family \cite{OpenAIChallenge2026,OpenAISolution2026}.

Several consequences follow from this syntax before any deep proof is
examined.  The family is chosen before the volume sequence, so this is
stronger than choosing a new finite model for each requested sequence.
The hierarchy is selected after refinement, so uniqueness across
refinements is not a conclusion.  All four fields of
$\operatorname{Properties}$ concern the same $T$, so the reconstruction,
operator-product estimates, and trace test cannot be supplied by unrelated
models.  These are elementary consequences of the order of quantifiers.
They are exactly the kind of details lost in a slogan such as ``a limit
exists.''

The public regularity predicate requires a positive logarithmic slope of
the bare action coefficient, uniformly in volume, continuous and
gauge-invariant lattice insertions, and one finite lattice stencil for each
token that works at every cutoff and volume.  The public
$\operatorname{Properties}$ predicate then bundles a Wilson-limit
statement, reconstructed fields, a separated-point operator-product
expansion, and a specified trace-based exclusion of free-field profiles
\cite{OpenAIChallenge2026}.  The strength of the result rests on the
definitions and proofs of these structures, not on the short final tactic
script that assembles them.

\section{A hierarchy with finite-degree local fields}

The formal hierarchy includes a Hilbert space, field labels, a Wightman
theory, finite-dimensional complex spaces $F_d$ of fields of bounded
engineering degree, maps from $F_d$ to the full label space, and a
surjective representation of those fields by lattice tokens.  It also
includes a compact space $X$ of phases with a probability measure, a
selected phase, correlation functions for the mixture and phases,
analytic continuations, OPE channels and coefficients, and a centered
degree-four trace field \cite{OpenAIChallenge2026}.

This is a structural resemblance to filtered algebra that can be used
carefully.  If $F_d$ is finite dimensional and increasing with $d$, then
one can choose successive complements and index finitely many new channels
at each degree.  Such choices support finite truncations of an expansion.
They do not imply that multiplication of quantum fields at the same point
is an everywhere-defined map $F_d\otimes F_e\to F_{d+e}$.
Coincidence limits need renormalization, and the OPE is an asymptotic
statement about separated fields.  Replacing that analytic statement by
an exact multiplication table would change the theorem.

Similarly, a quotient by correlation-null labels is a mathematically
natural way to obtain faithful fields.  If a preliminary label space $V$
maps linearly to correlation functions by $R$, the quotient
$V/\ker R$ injects into the image of $R$ by the first isomorphism
theorem.  This elementary fact explains why faithfulness of a quotient is
plausible; it does not prove that the chosen preliminary labels have the
desired Yang--Mills content.  That latter assertion requires comparison
with curvature polynomials and control of their null relations.  The
independent review consequently includes a separate operator-content
check instead of treating quotient faithfulness as a complete answer
\cite{YMAcceptance2026}.

\section{The two limits and the physical normalization}

Let $a_N\downarrow0$ denote the chosen spacings and $L_p\to\infty$
the physical side lengths.  A schematic smeared lattice correlation is
\begin{equation}
 C_{N,p}(v;f)=a_N^{4n}
 \sum_{x_1,\ldots,x_n}
 f(a_Nx_1,\ldots,a_Nx_n)
 \int\prod_{j=1}^n I_{N,p}(v_j,x_j;U)\,d\mu_{N,p}(U).
 \label{eq:openai-correlations}
\end{equation}
The original Wilson specification fixes the finite law, positions,
insertions, and separated-test convergence; the proof retains a full
ultraviolet limit for the chosen spacing sequence at fixed volume followed
by the asserted refined volume limit \cite{OpenAIWilson2026}.
The ordered limit is
$\lim_{k\to\infty}\lim_{N\to\infty}C_{N,p(q(k))}$.
No interchange of these limits is implicit.

The distinction can be understood without field theory.  The numbers
$f_{N,p}=N/(N+p)$ satisfy
$\lim_p\lim_N f_{N,p}=1$ and
$\lim_N\lim_p f_{N,p}=0$.  Therefore an ordered limit theorem cannot
be turned into a joint limit theorem without uniformity of an appropriate
kind.  In the gauge problem, locality, renormalization, and positivity make
the situation more structured, but the logical requirement remains.

The stronger acceptance specification also fixes a reference spacing and
a common reference coupling after renormalization-group evolution, and
requires the insertions to read out the same reference labels across
cutoffs and volumes.  Its comparison with another finite lattice
formalization includes the conversion of coupling conventions by the
matrix representation dimension \cite{YMAcceptance2026}.  These clauses
prevent a change in units or a rescaling of observables from masquerading
as a physically consistent limit.

In particular, the declaration named
\texttt{TwoPointScale.MassLimit} in the combined specification concerns
a positive limiting two-point calibration.  Its name must not be read as
the spectral condition \eqref{eq:mass-gap}.  A nonzero correlation
provides a scale-sensitive observable; the spectrum could still have
arbitrarily small positive energies.  The countermodels in the preceding
chapter demonstrate the logical difference explicitly.

\section{Reconstruction, phases, and test spaces}

The reconstructed theory includes a strongly continuous unitary
translation representation with spectrum in the forward cone, a dense
common domain for smeared fields, a normalized cyclic vacuum, adjoint and
regularity properties, Poincare covariance, and spacelike commutativity.
The selected phase has vacuum uniqueness.  The formal public definitions
specify these properties in the field system, rather than leaving the
meaning of ``Wightman theory'' as an informal comment
\cite{OpenAITheory2026}.  In this project these are statements of an
external formal development; our elementary semigroup arguments use their
consequences conditionally.

There are three separate classes of tests to keep apart.  First,
lattice-limit statements use smooth compactly supported tests away from
coincidences.  Second, the review's independent scalar OS comparison uses
Schwartz tests flat to every order on coincidence diagonals.  Third, the
reconstructed Wightman distributions act on unrestricted Schwartz tests.
The review compares the selected centered trace hierarchy to a separately
specified OS interface and growth condition
\cite{YMScalarOS2026}.  These domains have different powers to detect
contact terms, as the delta-distribution example above shows.

The mixed infinite-volume correlation table is related to phase
correlations by integration over $X$.  Each phase's reconstructed vacuum
uniqueness is compatible with nonuniqueness of the mixture's vacuum and
with disagreement between different infinite-volume subsequences.  The
combined acceptance theorem keeps the selected scalar field, phase, trace,
and translation tensor attached to one construction witness
\cite{YMCombined2026}.  This coherence is mathematically significant.
It does not remove the remaining uniqueness questions by definition.

\section{The trace law, non-freeness, and OPE estimates}

For the normalized centered stress trace $\Theta$ in the selected phase,
the public theorem includes a positive leading short-distance constant:
\begin{equation}
 \lim_{r\downarrow0}
 r^8\log^2(1/r)
 S_{\Theta,\Theta}(-re_0/2,re_0/2)
 =\frac{3\dim_{\mathbb R}\mathfrak g}{\pi^4}.
 \label{eq:openai-trace}
\end{equation}
The independent trace adapter expresses the same limit in all spatial
directions for the same centered field and normalization
\cite{YMLeadingTrace2026}.  The detailed coefficient is convention
dependent; copying it while changing the normalization of the stress
tensor would be incorrect.

The public non-freeness test excludes finite combinations of specified
mass-shell Wick profiles, including finite products of kernels with
nonnegative masses and homogeneous numerators.  This is a precise
exclusion class, and the formal target describes it explicitly
\cite{OpenAIChallenge2026}.  It should not be expanded into an assertion
excluding every possible Gaussian representation, every generalized free
field with a continuous mass spectrum, or every conceivable equivalence
of theories.  These stronger statements would need new definitions and
proofs.  Nor does \eqref{eq:openai-trace} establish the gap: multiplying
the examples in \eqref{eq:spectral-countermodel} by the appropriate
positive constant gives exactly the same leading coefficient with
different infrared supports.

The separated OPE has an input cluster of fields $v_i$ of total degree
$w=\sum_i d_i$, a cluster radius $r$, and spectators kept a fixed
positive distance away.  At cutoff degree $D$ it has the form
\begin{equation}
 S^s(v(rx),u(y))=
 \sum_{d_\alpha\leq D}c_\alpha(v;rx)S^s(o_\alpha(0),u(y))
 +R_{D,s}(r),
 \label{eq:openai-ope}
\end{equation}
with bounds of the type
$|c_\alpha|\leq r^{d_\alpha-w-\varepsilon}$ and
$|R_{D,s}(r)|\leq r^{D+1-w-\varepsilon}$ for all sufficiently small
$r$, uniformly over the prescribed separated configurations and phases.
The same coefficients serve all spectators, phases, and truncation
degrees.  The independent OPE check restates such bounds directly and
derives them from the source theorem \cite{YMOPE2026}.

Here ``all orders'' means arbitrarily high \emph{field-degree truncation}
in this expansion.  It is not a claim that every perturbative loop
coefficient has been compared with a particular renormalization scheme.
The difference is visible in the variables: $D$ indexes the degrees of
local fields, whereas loop order indexes powers of a coupling in a
perturbative series.  An equality between the two kinds of information
would be another theorem.  Maintaining this distinction makes the OPE
result more useful, because it specifies the exact remainder one can
legitimately use in further arguments.

\section{A stronger acceptance target and the scope of checking}

The original short public target did not retain every connection between
the lattice construction and the continuum witness.  The review's
\texttt{CombinedAcceptance} specification explicitly keeps one family,
one extracted table, and one final construction object with those
certificates.  Its proof applies the stronger source construction theorem
once and assembles the additional consequences for that same object
\cite{YMAcceptance2026,YMCombined2026}.  This is a concrete example of why
verifying a theorem's proof term and verifying that its type expresses the
intended scientific claim are distinct activities.

The recorded independent check reports a clean build of 9,470 project
modules using Lean 4.32.0, with 9,477 supplied source files matched to a
manifest.  It reports no project axioms, no direct placeholder dependency,
and no unsafe declarations in the inventoried project declarations.  A
reference challenge contains a placeholder, but the solution does not
import that reference proof.  The reported principal proof dependencies
use the standard axioms of propositional extensionality, classical choice,
and quotient soundness \cite{YMBuildEvidence2026}.

The group's separate Comparator records cover the public endpoint, the
combined acceptance endpoint, and the OPE endpoint.  They include controls
that reject a changed target, an added forbidden axiom, and a changed
definition.  The historical macOS launcher did not provide Linux process
isolation, and the check used Comparator's built-in Lean kernel rather
than another independently implemented kernel
\cite{YMComparatorEvidence2026}.  These qualifications describe the
actual experiment.  They do not negate a kernel check, and they should
not disappear when the check is summarized.

This project distinguishes four levels of confidence.  A well-typed
declaration verifies a particular proposition relative to its dependencies.
An axiom audit identifies the assumptions used.  A specification comparison
checks that a fixed target was not replaced.  Semantic review assesses
whether the fixed target captures the mathematical and physical claim.
The last step cannot be reduced to searching for the word
\texttt{sorry}: a theorem with a vacuous hypothesis can have a completely
valid proof and still express the wrong scientific question.  Conversely,
a separate unfinished reference stub is not an assumption of a proof that
never imports it.  Both observations follow from the dependency structure
of formal logic.

\section{What can be built on the result now}

The most useful next contributions are statements whose inputs can be
bound to the already constructed objects.  Examples include a precise
uniqueness criterion for the phase mixture, a criterion that identifies
different volume subsequences, or a volume-uniform decay estimate on a
total observable domain.  Each should specify the family, normalization,
test spaces, and phase before stating a conclusion.  The spectral
proposition in the preceding chapter explains exactly how a successful
uniform estimate would yield a gap.

An algebraic improvement could be valuable without resolving the infrared
problem.  A clearer description of the filtered local-field spaces might
simplify the verification of null relations.  A finite presentation of
low-degree relations could make the operator-content interface easier to
audit.  A normalized family of finite complexes could organize Ward
identities and cohomological constraints.  Such work would need to show
that its algebraic relations are respected by the same renormalized
correlation hierarchy, rather than merely by a classically related
polynomial ring.

The appropriate positive conclusion is therefore substantial and precise:
the current public development supplies a formally verified external
construction with documented independent checks, and it supplies concrete
objects to which further theorems can refer.  The remaining gap is not an
absence of numerical plausibility.  It is the absence, in the reviewed
claim, of a theorem giving a strictly positive lower bound on the
nonvacuum energy spectrum in the infinite-volume continuum theory.  The
present manuscript does not supply that missing theorem and does not label
its classical or finite-dimensional propositions as substitutes for it.

\chapter{An explicit arithmetic--geometric bridge through elliptic tori}
\label{chap:elliptic-gauge-bridge}

\section{Orders acting on a complex elliptic curve}

A concrete example can connect orders, elliptic curves, unitary gauge
fields, and elliptic operators without making a conjectural identification.
Let $\Lambda=\mathbb Z\omega_1+\mathbb Z\omega_2\subset\mathbb C$
with $\operatorname{Im}(\omega_2/\omega_1)>0$, and form the complex
torus $E=\mathbb C/\Lambda$.  Its algebraic realization as a complex
elliptic curve is the usual uniformization theorem.  The elementary
constructions below begin with this displayed torus, so they do not need
uniformization as an unproved step.

For $\alpha\in\mathbb C$, multiplication on $\mathbb C$ descends to
a map $[z]\mapsto[\alpha z]$ on $E$ exactly when
$\alpha\Lambda\subseteq\Lambda$.  Necessity follows because each
$\lambda\in\Lambda$ represents the same point as zero.  Sufficiency
follows because $z-z'\in\Lambda$ then implies
$\alpha z-\alpha z'\in\Lambda$.  Thus
\[
 \operatorname{End}_{\mathrm{mult}}(E)
 =\{\alpha\in\mathbb C:\alpha\Lambda\subseteq\Lambda\}
\]
is a ring.  If it contains a nonreal $\alpha$, it lies in an imaginary
quadratic number field and is an order in that field.  To see the essential
algebra, represent multiplication by $\alpha$ on the chosen lattice basis
by a matrix $M_\alpha\in M_2(\mathbb Z)$.  Its characteristic
polynomial is monic of degree two, annihilates $\alpha$, and has negative
discriminant when $\alpha$ is nonreal.  The finite-index assertion for the
endomorphism ring is the standard complex-multiplication structure theorem;
the matrix calculation identifies its integrality mechanism.  One can make
the order assertion explicit.  Rescale so that
$\Lambda=\mathbb Z+\mathbb Z\tau$.  Write $\alpha=a+c\tau$ and
$\alpha\tau=b+d\tau$ with integers $a,b,c,d$.  Since $\alpha$ is
nonreal, $c\ne0$, and $c\tau^2+(a-d)\tau-b=0$.  Thus
$K=\mathbb Q(\tau)$ is imaginary quadratic.  Every other multiplier
$\beta$ has $\beta\cdot1\in\Lambda$, so $\beta\in K$.
Its integral matrix on $\Lambda$ makes it an algebraic integer, hence
the multiplier ring lies in $\mathcal O_K$.  It contains
$\mathbb Z[\alpha]$, a rank-two sublattice of $\mathcal O_K$.
Therefore it too is a rank-two subring of finite index in
$\mathcal O_K$, which is exactly an order.  Here the usual fact that
$\mathcal O_K$ is a rank-two lattice is the number-field theorem
established or cited in the arithmetic chapters.

This is an actual use of an order acting on a lattice.  It concerns
endomorphisms of a complex torus and should be compared with the arithmetic
endomorphism-order results in the earlier chapters.  It does not replace
the separate problem of computing endomorphisms or discrete logarithms
over a finite field.

\section{The same lattice controls flat unitary line bundles}

A unitary character $\chi:\Lambda\to U(1)$ determines a line bundle
\[
 L_\chi=(\mathbb C\times\mathbb C)/\sim,
 \qquad (z,v)\sim(z+\lambda,\chi(\lambda)v).
\]
Here $\chi$ is the automorphy character: with this quotient convention, parallel-transport holonomy is $\chi^{-1}$. The trivial connection on the product descends because the multipliers
are constant.  It is flat and unitary.  Conversely, parallel transport of
a flat unitary line bundle around the two lattice generators gives its
holonomy character.  In rank one, conjugacy of holonomy characters is
trivial.  The general relation between stable bundles and unitary
representations is a deep theorem of Narasimhan and Seshadri; Atiyah's
classification gives the special structure of bundles on elliptic curves
\cite{YMNS1965,YMAtiyah1957}.

For this explicit rank-one model, all characters are
\begin{equation}
 \chi_{s,t}(m\omega_1+n\omega_2)
 =\exp(2\pi i(ms+nt)),\qquad
 (s,t)\in(\mathbb R/\mathbb Z)^2.
 \label{eq:torus-character}
\end{equation}
Indeed a homomorphism out of a free abelian group is uniquely determined
by the images of its basis, and each unit complex number has an argument
modulo $2\pi$.  Tensor product of line bundles corresponds to multiplying
characters, hence adding $(s,t)$.  The characters with $\chi^n=1$ are
exactly the $n^2$ points with $s,t\in\frac1n\mathbb Z/\mathbb Z$.
This explains how finite torsion and finite Fourier characters occur inside
a moduli space of flat connections.

\begin{proposition}[Order action and dual action]
If $\alpha\Lambda\subseteq\Lambda$, pullback by
$[z]\mapsto[\alpha z]$ sends $L_\chi$ to the flat line bundle with
character $\chi\circ\alpha$.  If $M_\alpha$ represents multiplication
on column lattice coordinates, the corresponding character parameters
transform by $M_\alpha^{\mathsf T}$ modulo $\mathbb Z^2$.
\end{proposition}

\begin{proof}
Under pullback, the constant multiplier for a lattice translation $\lambda$ is the multiplier for $\alpha\lambda$, namely $\chi(\alpha\lambda)$.  If $x=(m,n)^{\mathsf T}$ denotes lattice
coordinates and $u=(s,t)^{\mathsf T}$ denotes character coordinates,
then \eqref{eq:torus-character} is $\exp(2\pi i u^{\mathsf T}x)$.
Precomposition replaces $x$ by $M_\alpha x$, giving
$u^{\mathsf T}M_\alpha x=(M_\alpha^{\mathsf T}u)^{\mathsf T}x$.
\end{proof}

The transpose is the same algebraic operation that appears in dual
lattices and Fourier transforms.  What unifies these examples is the
evaluation pairing between a module and its characters.  That pairing can
be formalized directly.  It should not be enlarged into a claim that
arithmetic elliptic curves, all gauge moduli spaces, and quantum hidden
subgroup problems are equivalent categories.

\section{An elliptic operator with a completely explicit spectrum}

Take a rectangular torus with side lengths $L_1,L_2>0$ and automorphy parameters
$(s,t)$.  Sections of its flat line bundle can be represented by smooth
functions satisfying
\[
 f(x+L_1,y)=e^{2\pi is}f(x,y),\qquad
 f(x,y+L_2)=e^{2\pi it}f(x,y).
\]
The nonnegative covariant Laplacian is represented by
$-\partial_x^2-\partial_y^2$ with these boundary conditions.  Its
Fourier modes are
\[
 f_{m,n}(x,y)=
 \exp\left(2\pi i\left(\frac{m+s}{L_1}x+
                              \frac{n+t}{L_2}y\right)\right),
 \qquad (m,n)\in\mathbb Z^2,
\]
with eigenvalues
\begin{equation}
 \lambda_{m,n}(s,t)=4\pi^2
 \left(\frac{(m+s)^2}{L_1^2}+\frac{(n+t)^2}{L_2^2}\right).
 \label{eq:twisted-spectrum}
\end{equation}
The boundary conditions and eigenvalue equation follow by direct
substitution.  Completeness follows by multiplying by
$\exp(-2\pi i(sx/L_1+ty/L_2))$, which turns such sections into
periodic functions, and applying the usual complete Fourier basis of the
rectangular torus.

There is a zero mode exactly for the trivial character: a zero eigenvalue
requires $m+s=n+t=0$ for integers $m,n$.  For every nontrivial character
the lowest eigenvalue is positive, but it tends to zero as the character
approaches the trivial one.  This example simultaneously exhibits a
moduli parameter, a jump in a kernel dimension, and failure of a uniform
spectral bound over moduli.  It can be plotted exactly, without numerically
solving a differential equation.  A diagram should label its two displayed
automorphy coordinates and should describe a four-dimensional parameter
view as a projection or slice if side lengths are varied too.

The corresponding Dolbeault operator
$\overline\partial_{L_\chi}:\Omega^{0,0}(L_\chi)\to
\Omega^{0,1}(L_\chi)$ is an elliptic complex on the complex curve.
Riemann--Roch gives index zero for a degree-zero line bundle on a
genus-one curve.  Thus the index remains constant while individual kernel
dimensions and positive eigenvalues vary.  The explicit flat spectrum
already demonstrates the crucial distinction.  Atiyah--Singer index
theory organizes stable topological information; it does not fix every
eigenvalue in a family.

\section{How this bridge relates to cryptography}

The finite $n$-torsion characters above form
$(\mathbb Z/n\mathbb Z)^2$, and their evaluation pairing is the
familiar finite Fourier pairing.  A quantum Fourier transform can sample
characters of a finite abelian group when that group is presented with
efficient reversible operations.  The mathematical resemblance is exact
at the level of the pairing.  The complexity statement needs more data:
an encoding of elements, circuits for the group law, access to the hidden
function, and a cost model.  A geometric interpretation of characters
does not automatically provide those resources.

For cryptographic elliptic curves over $\mathbb F_q$, the public
computational problem usually concerns a specified finite subgroup of
$E(\mathbb F_q)$ and scalar multiplication in that subgroup.  The
complex torus model instead uses a lattice in $\mathbb C$.  Relating the
two requires arithmetic reduction or lifting with its own hypotheses,
and a proposed algorithm must preserve the relevant unknown and control
the cost of every conversion.  No such efficient conversion is supplied
by the flat-line-bundle example.  Its contribution is a rigorous
conceptual connection, useful for organizing structures and formal
interfaces, rather than a claimed cryptanalytic improvement.

There is a parallel lesson for lattice-based cryptography.  A module over
an order can be equipped with a positive form and embedded as a Euclidean
lattice.  Its dual then represents characters, and Fourier methods can
be expressed geometrically.  Hardness, however, concerns quantitative
approximation or average-case distributions.  The purely algebraic
existence of a dual or a character basis does not bound the complexity of
finding a shortest vector, recovering a noisy secret, or implementing the
needed state preparation.  Those are additional statements and belong
in the algorithmic part of the research program.

\section{A precise synthesis and a testable research program}

The preceding constructions establish a network of actual mathematical
maps.  Orders act on integral lattices; a complex multiplication order acts
on an elliptic torus; pullback acts contragrediently on its unitary
characters; characters specify flat line bundles; their covariant
differentials form complexes; elliptic symbols yield index classes; and
finite gauge discretizations provide a setting for quantum reconstruction
when the necessary analytic limits are proved.  Each arrow has a defined
domain and codomain.  A useful new result must strengthen one of those
arrows, or supply a missing quantitative estimate, rather than merely
draw the same objects closer together in a diagram.

Three research directions are consequently concrete.  First, formalize the
elementary matrix, character, and deformation identities and connect them
to existing verified libraries through explicit equivalences.  Second,
develop finite algebraic presentations of selected gauge-observable
relations and check their compatibility with one pinned continuum
construction.  Third, investigate quantitative inequalities on a total
observable domain, retaining dependence on cutoff, volume, normalization,
and phase.  Success in the first direction provides reliable infrastructure;
success in the second may improve exposition and auditing of the
construction; a sufficiently strong success in the third could address the
remaining mass-gap question.

These directions have different completion criteria.  A matrix identity
is complete when its stated assumptions yield its conclusion in the
kernel.  A compatibility theorem is complete when the diagram of maps
commutes for the actual objects used by the external construction.
A mass-gap argument is complete only when its uniform estimates and
limit passages prove \eqref{eq:mass-gap} for the required nontrivial
quantum Yang--Mills theory.  Keeping those criteria explicit allows
substantial intermediate progress to be recognized without relabeling
it as a solution of an unresolved problem.
